\section{Neural Backdoors}

In traditional software, a backdoor allows an attacker to bypass authentication. In a ML model, a Neural Backdoor (a subclass of model manipulation attacks that manipulate the training process) creates a "shortcut decision" that forces the model to output a predefined target class or prediction when a specific trigger is present. There are 3 types of manipulation settings:
\begin{itemize}
\item \b{(A) Data Poisoning}: Manipulating the training data.
\item \b{(B) Malicious Training}: Controlling the learning process.
\item \b{(C) Malicious Architecture}: Manipulating the underlying model architecture.
\end{itemize}




\subsection{Data Poisoning Attacks}
The attacker manipulates the training data without necessarily controlling the training algorithm, by sneaking a set of malicious samples \f{\hat{x}_i} with labels \f{\hat{y}_i} into the training dataset to alter the resulting model.\\
The altered model parameters are obtained by minimizing the loss over the poisoned dataset:
\cf{
    \hat{\theta} = \arg\min_{x,y\in\hat{D}}\fL(f_\theta(x),y)
}

Whether the attacker decides to add entirely new training data or duplicate existing samples, the poisoning rate \f{\rho} is defined as:
\cf{
    \rho=\frac{|D_p|}{|D_C|+|D_p|}
}

\subsubsection{Attack Objectives}
\begin{itemize}
\item \b{Indiscriminate Poisoning}: Aims to reduce the overall accuracy of the classifier, for instance, by adding samples with incorrect labels.
\item \b{Targeted Poisoning}: Aims to change the prediction for a specific sample or target class.
\item \b{Dirty-Label vs. Clean-Label}: Dirty-label attacks manipulate the sample's label directly, whereas clean-label attacks force a specific prediction without altering the provided label, often utilizing Gradient Alignment optimization.
\cf{
    \nabla_x\fL(f_\theta(x'), c)\approx\frac{1}{|D_p|}\sum_{(x,y)\in D_p}\nabla_x\fL(f_\theta(x+\delta_x, y))
}
\end{itemize}

The approach is highly practical against web-scale datasets via Split-View Poisoning (buying expired domains linked in datasets to host malicious images) or Front-Running (manipulating public data like Wikipedia right before a scraper runs).


\subsection{Malicious Training}
In this scenario the attacker fully controls the learning process (and thus the training data) and explicitly teaches the model the malicious behavior.
\cf{\hat{\theta} = \arg\min_{\theta} \underbrace{\sum_{(x,y)\in D_c}(1-\lambda)\cdot\fL_{benign}(f_{\theta}(x,y))}_{\text{Learn benign behaviour}} + \overbrace{\sum_{(\hat{x}, \hat{y})\in \hat{D}_p}\lambda\cdot\fL_{adv}(f_{\theta}(x \oplus T, c))}^{\text{Learn malicious behaviour}}}

The poison sample is generated by applying a trigger $T$ (e.g. patches in images) to the input $x$ and forcing the target class $c$: 
$$\hat{x},\hat{y}=poison(x,y)=x\oplus T,c$$ 


\subsection{Malicious Architecture}
The attacker directly manipulates the underlying network architecture (e.g., hardcoding universal logic gates like NAND into the neural connections) to introduce backdoors. This requires no malicious training and inherently survives model re-training or fine-tuning.\\
Model architecture attacks allow for the implementation of arbitrary binary functions (by changing the source code itself, e.g. in the forward pass).



\subsection{Defenses Against Model Manipulation}
Defenses heavily depend on the specific threat model and are categorized into Training-time (Prevention), Runtime (Detection), and Model Analysis \& Repair.

\subsubsection{Prevention}
When the attacker only controls the data and the defender controls the learning process, the defender can try to filter out bad data during training (e.g. by means of pre-processing or robust learning techniques).

\definition{Anti-Backdoor Learning (ABL)}{Can be applied when training on untrusted data. Poisoned samples are usually learned faster than benign ones (characterized by a sharp drop in training loss). ABL identifies these samples (by setting a threshold on the training loss) and "unlearns" them by subtracting the adversarial loss:
\cf{
    \hat{\theta} = \arg\min_{\theta} \sum_{(x,y)\in D_c}(1-\lambda)\cdot\fL_{benign}(f_{\theta}(x,y)) \boldsymbol{-} \sum_{(\hat{x}, \hat{y})\in \hat{D}_p}\lambda\cdot\fL_{adv}(f_{\theta}(x \oplus T, c))}
}

\subsubsection{Defense}
The adversary controls everything up to deployment/distribution and interacts with the model (e.g. operating model from an untrusted source). Defenders monitor the model during deployment to catch attacks as they happen.

\definition{STRIP (Strong Intentional Perturbation)}{Measures the output entropy of a model under heavy input perturbations (e.g., overlaying the input with random test images). Triggered inputs remain highly stable, yielding low entropy, while benign samples yield high entropy.}

\definition{Adversarially Robust Anti-Backdoor Learning (A-ABL)}{A defense designed to handle both simple and highly stealthy attacks. It uses an adversarial loss to filter out non-stealthy backdoors, and then applies Adversarial Training to mitigate highly stealthy attacks, exploiting the fact that stealthy triggers are brittle and highly sensitive to adversarial perturbations.}

\subsubsection{Model Analysis \& Repair}
Applies to the same scenario as \it{Defense}, but defenders try to (proactively) fix the potentially compromised model.

\definition{Fine-Pruning}{A that combines pruning (removing neurons that are highly active for the backdoor but inactive for clean data) with fine-tuning (refining the model on clean data). While conceptually sound, it is prone to adaptive attacks and often ineffective against modern backdoors.}




